% Generated draft: requires mathematical and source review.
\documentclass[11pt]{article}
\usepackage[margin=25mm]{geometry}
\usepackage{fontspec}
\setmainfont{Latin Modern Roman}
\usepackage{amsmath,amssymb,mathtools,unicode-math}
\setmathfont{Latin Modern Math}
\usepackage{enumitem,needspace,xurl,hyperref}
\setlength{\parindent}{0pt}
\setlength{\parskip}{5pt}
\setlength{\emergencystretch}{3em}
\newcommand{\N}{\mathbb{N}}
\newcommand{\Z}{\mathbb{Z}}
\newcommand{\Q}{\mathbb{Q}}
\newcommand{\R}{\mathbb{R}}
\newcommand{\C}{\mathbb{C}}
\newcommand{\F}{\mathbb{F}}
\newcommand{\A}{\mathbb{A}}
\newcommand{\PP}{\mathbb{P}}
\newcommand{\E}{\mathbb{E}}
\newcommand{\eps}{\varepsilon}
\newcommand{\dd}{\,\mathrm{d}}
\begin{document}
\section*{UnsolvedMath ID: 20000228}
\textbf{Source problem code:} catalogue entry
\subsection{Definitions and mathematical statement}
Let $X_s\to\PP^2_{\C}$ be the blowup of $s$ very general points, let $L$ be the pullback of a line, and let $E_i$ be the exceptional divisors. For a divisor $D=dL-\sum_{i=1}^s m_iE_i$, the AIM question asks for a convex-method proof that $D$ is not big whenever the intended Nagata prediction gives $h^0(X_s,\mathcal O(D))=0$. Big means positive volume, or equivalently interior membership in the pseudoeffective cone for numerical divisor classes. The source does not specify the precise vanishing prediction or whether equal or arbitrary multiplicities are intended. Its remark says such a result might imply a strong form of Nagata's conjecture; this is a stated motivation, not an established implication here.

The supplied report takes nonnegative multiplicities and studies a particular toolbox of nef square-subset classes
\[
 A_{n,S}=nL-\sum_{i\in S}E_i,\qquad n\geq1,\quad |S|=n^2.
\]
For fixed $n^2\leq s$, it describes their nonnegative span as the homogenized hypersimplex. Explicitly a class $aL-\sum b_iE_i$ is in that cone precisely under the corresponding inequalities and relation
\[
 a\geq0,\qquad 0\leq b_i\leq a/n,\qquad\sum_i b_i=na.
\]
Write $M_q$ for the sum of the $q$ largest $m_i$. The claimed sharp fixed-$n$ separating condition is
\[
 M_{n^2}\geq nd.
\]
The relevant intersection is $D\cdot A_{n,S}=nd-\sum_{i\in S}m_i$.

Allowing all square sizes and arbitrary further very general blown-up points with multiplicity zero, the report gives the exact threshold, for $s\geq1$,
\[
 \beta_{\mathrm{square}}=
 \max\left\{\max_{1\leq n\leq\lfloor\sqrt s\rfloor}\frac{M_{n^2}}n,
 \frac{\sum_{i=1}^s m_i}{\lceil\sqrt s\rceil}\right\}.
\]
It claims that $d\leq\beta_{\mathrm{square}}$ certifies non-bigness, $d<\beta_{\mathrm{square}}$ certifies non-pseudoeffectivity, and $d>\beta_{\mathrm{square}}$ cannot be certified by any nonnegative combination of these square rays, even after the permitted auxiliary extensions. This last statement limits this toolbox; it does not prove bigness or pseudoeffectivity above the threshold.

The report further states that homogeneous weights are the sharp worst case at fixed total multiplicity, that the resulting approximation factor is at least $\sqrt3/2$ for $s\geq10$, and that an explicit integer rounding gap is available. The benchmark and rounding formula are not supplied in the record and require clarification. All claimed results are partial; in particular the report expressly does not solve nonsquare Nagata.
\subsection{Short English statement}
Find convex certificates for non-bigness in the range where Nagata predicts no sections on a very general plane blowup. The supplied report computes the exact power and limitation of certificates built from nef square-subset rays, including auxiliary zero-weight points. Its threshold distinguishes non-bigness from non-pseudoeffectivity and does not settle nonsquare Nagata.
\subsection{Sources}
\begin{itemize}
\item Local registry: \url{data/unsolvedmath/problems.json}, permanent ID $20000228$, source code \texttt{catalogue entry}; original question, Remark 7.2, and supplied research summary.
\item AIM workshop ``Degenerations in algebraic geometry'', problem-session notes dated September 7, 2016, Problem 7.1: \url{https://aimath.org/pastworkshops/degenalggeomproblems.pdf}. Canonical location \texttt{aim-algebraic-geometry-notes.json}, \texttt{notes[227]}; original literature field empty. No PDF was opened here.
\item Supplied machine-generated research attempt 1: partial result, problem open at the run and currently partially solved; rigor metadata ``proved'', novelty confidence low. No full authoritative citation or detailed proof for the square-ray threshold and approximation assertions is supplied. This non-peer-reviewed aid requires independent verification before citation.
\end{itemize}
\end{document}
